\chapter*{LỜI CẢM ƠN}
\addcontentsline{toc}{chapter}{LỜI CẢM ƠN}
\fontsize{13}{15}\selectfont
Trong suốt quá trình thực hiện đồ án tốt nghiệp, em đã nhận được sự quan tâm, hướng dẫn và hỗ trợ tận tình từ nhiều cá nhân và tập thể. Nhân dịp này, em xin bày tỏ lòng biết ơn chân thành và sâu sắc nhất.

Trước tiên, em xin gửi lời cảm ơn trân trọng đến TS. Nguyễn Kiêm Hùng — cán bộ hướng dẫn chính — và TS. Mai Linh — cán bộ đồng hướng dẫn — đã tận tình định hướng đề tài, dành thời gian trao đổi chuyên môn và đưa ra những nhận xét xác đáng trong suốt quá trình nghiên cứu. Những chỉ dẫn của thầy cô về xử lý tín hiệu số, tối ưu hóa mô hình trên nền tảng TinyML và quản lý tài nguyên bộ nhớ trên vi điều khiển là nền tảng quan trọng giúp em hoàn thành đồ án này.

Em xin trân trọng cảm ơn toàn thể quý thầy cô giảng viên thuộc Khoa Điện tử Viễn thông, Ngành Kỹ thuật Máy tính, Trường Đại học Công nghệ – Đại học Quốc gia Hà Nội đã truyền đạt nền tảng kiến thức vững chắc trong suốt bốn năm đào tạo. Những kiến thức về kiến trúc máy tính, hệ điều hành thời gian thực và trí tuệ nhân tạo chính là cơ sở lý luận và thực tiễn không thể thiếu để triển khai đề tài này.

Cuối cùng, em xin gửi lời tri ân sâu sắc đến gia đình và bạn bè đã luôn đồng hành, động viên và tạo điều kiện thuận lợi trong suốt thời gian thực hiện đồ án.

Do giới hạn về thời gian và kinh nghiệm thực tiễn, đồ án không tránh khỏi những thiếu sót nhất định. em kính mong nhận được sự góp ý, chỉ bảo của quý thầy cô trong Hội đồng chấm đồ án để công trình được hoàn thiện hơn.

Em xin chân thành cảm ơn!

\vspace{0.5cm}
\begin{flushright}
    Quách Ngọc Quang\hspace*{1.3cm}
\end{flushright}

\chapter*{LỜI CAM ĐOAN}
\addcontentsline{toc}{chapter}{LỜI CAM ĐOAN}
\fontsize{13}{15}\selectfont
Tôi tên là Quách Ngọc Quang, sinh viên năm tư ngành Kỹ thuật máy tính, Trường Đại học Công nghệ – Đại học Quốc gia Hà Nội. Tôi xin cam đoan đồ án tốt nghiệp với đề tài \textit{"Thiết kế hệ thống AIoT phát hiện và cảnh báo bất thường cho máy giặt dân dụng sử dụng Autoencoder không giám sát trên thiết bị biên"} là công trình nghiên cứu của riêng tôi, được thực hiện dưới sự hướng dẫn khoa học của TS. Nguyễn Kiêm Hùng và TS. Mai Linh.

Các số liệu, kết quả thực nghiệm trình bày trong đồ án về độ trễ suy luận, sai số tái tạo (MAE) và các chỉ số lượng tử hóa mô hình là hoàn toàn trung thực và chưa từng được công bố trong bất kỳ công trình nào khác. Hệ thống phần cứng dựa trên vi điều khiển Seeed Studio XIAO ESP32-S3, mã nguồn firmware vận hành đa lõi trên FreeRTOS và ứng dụng giám sát đi kèm đều do tôi tự nghiên cứu, phát triển và tối ưu hóa. Các tài liệu tham khảo và mã nguồn mở sử dụng trong dự án đều được trích dẫn nguồn gốc rõ ràng.

Nếu có bất kỳ sự gian lận hoặc vi phạm quy chế nào về tính nguyên bản của công trình, tôi xin chịu hoàn toàn trách nhiệm trước Hội đồng bảo vệ và Ban Giám hiệu nhà trường.

\vspace{0.5cm}

\begin{flushright}
    Hà Nội, ngày ... tháng ... năm 2026\par
    Tác giả\hspace*{2.8cm}\par
    \vspace{3.5cm}
    Quách Ngọc Quang\hspace*{1.3cm}
\end{flushright}

\chapter*{PHÊ DUYỆT CỦA CÁN BỘ HƯỚNG DẪN}
\addcontentsline{toc}{chapter}{PHÊ DUYỆT CỦA CÁN BỘ HƯỚNG DẪN}
\fontsize{13}{15}\selectfont

...

\vspace{1.2cm}

\vfill

\begin{flushright}
    Hà Nội, ngày ... tháng ... năm 2026
\end{flushright}

\vspace{0.6cm}

\noindent
\begin{tabularx}{\textwidth}{>{\centering\arraybackslash}X>{\centering\arraybackslash}X}
    \textbf{Cán bộ hướng dẫn} & \textbf{Cán bộ đồng hướng dẫn} \\[3.4cm]
    \textbf{TS. Nguyễn Kiêm Hùng} & \textbf{TS. Mai Linh}
\end{tabularx}

\vspace{0.8cm}
